\section{Related work}

LegalBench established a broad expert-built benchmark for legal language
understanding \citep{guha2023legalbench}; CUAD supplied a contract-analysis
corpus with span and category supervision \citep{hendrycks2021cuad}. ContractNLI
connects natural-language clauses to entailment, contradiction, and
insufficient-information judgments \citep{koreeda2021contractnli}. These
benchmarks motivate separate evidence and behavior units, but they do not by
themselves specify a fail-closed compiler contract.

Formalization systems expose a complementary tension. Grammars of Formal
Uncertainty reports that autoformalization can improve or harm downstream
behavior depending on the grammar and task \citep{ganguly2025grammars}.
Catala demonstrates that legal rules can be represented in a domain-specific
language with executable semantics \citep{bavard2021catala}. RuleArena studies
rule-following evaluation under adversarial conditions \citep{wang2024rulearena}.
Recent work on legal-text decision models and policy guards further motivates
grounded evidence and explicit uncertainty \citep{graus2026legaltext,zhou2026lgmt}.

LEXEC differs in emphasis. It treats provenance, refusal, schema version, and
unrun status as part of the evaluation object, and it makes the boundary
between source fidelity and executable behavior operational. The resulting
instrument can host existing benchmarks without claiming that a compiler's
internal checks replace expert judgments.
